\documentclass[10pt]{extarticle}
\usepackage[a4paper,margin=1.5cm]{geometry}
\usepackage[T1]{fontenc}
\usepackage{amsmath,amssymb}
\usepackage{graphicx}
\usepackage[export]{adjustbox}
\usepackage{multicol}
\usepackage{enumitem}
\usepackage{booktabs,array,tabularx}
\setlength{\columnsep}{0.7cm}
\setlength{\parindent}{0pt}
\setlength{\parskip}{3pt}
\pagestyle{empty}
\makeatletter
\newcommand{\qfigdef}[4]{\@namedef{qf@file@#1}{#2}\@namedef{qf@w@#1}{#3}\@namedef{qf@wide@#1}{#4}}
\newcommand{\qf@check}[1]{\@ifundefined{qf@file@#1}{\PackageError{score2}{Unknown figure id #1}{}}{}}
\newcommand{\qf@img}[4][]{\includegraphics[width=#2\dimexpr\@nameuse{qf@w@#4}pt\relax,max width=#3,max height=0.45\textheight,#1]{\@nameuse{qf@file@#4}}}
\newif\ifqf@full
\newif\ifqf@boxed
\newsavebox{\qf@box}
\newcommand{\qf@open}{\global\qf@fulltrue\par\global\qf@boxedtrue
  \setbox\qf@box\vbox\bgroup\hsize\textwidth\linewidth\textwidth}
\newcommand{\qf@emit}[1]{\par\egroup\global\qf@boxedfalse
  \ifdim\dimexpr\ht\qf@box+\dp\qf@box\relax<#1\textheight\noindent\usebox{\qf@box}\par\else\unvbox\qf@box\fi}
\newcommand{\qfullsplit}{\ifqf@boxed\qf@emit{0.85}\fi}
\newcommand{\qf@close}{\ifqf@boxed\qf@emit{0.45}\fi\global\qf@fullfalse}
\newcommand{\qfullbegin}{\end{multicols}\qf@open}
\newcommand{\qfullend}{\qf@close\begin{multicols}{2}}
\newcommand{\qfullstart}{\qf@open}
\newcommand{\qfulldone}{\qf@close}
\newcommand{\qfig}[2][1]{\qf@check{#2}%
  \ifqf@full\qf@blk{#1}{#2}\else
  \ifnum\@nameuse{qf@wide@#2}=1 \expandafter\@firstoftwo\else\expandafter\@secondoftwo\fi
  {\end{multicols}\begin{center}\qf@img{#1}{\textwidth}{#2}\end{center}\begin{multicols}{2}}%
  {\qf@blk{#1}{#2}}\fi}
\newcommand{\qf@blk}[2]{\par\smallskip{\centering\qf@img{#1}{\linewidth}{#2}\par}\smallskip}
\newcommand{\qfigi}[2][1]{\qf@check{#2}\qf@img[valign=c]{#1}{\linewidth}{#2}}
\makeatother
\newcommand{\papersec}[1]{\par\bigskip{\centering\Large\bfseries PAPER-#1\par}\medskip}
\newcommand{\partsec}[2]{\par\medskip\noindent{\large\bfseries PART-#1: #2}\par\smallskip}
\newcommand{\sect}[1]{\par\smallskip\noindent\textbf{#1}\par}
\newcommand{\qn}[1]{\par\smallskip\noindent\textbf{#1.}~}
\newcommand{\paraq}[2]{\par\smallskip\noindent\textbf{Paragraph for Questions #1 and #2}\par}
\newenvironment{qoptions}{\begin{itemize}[leftmargin=1.8em,itemsep=1pt,topsep=2pt,label={}]}{\end{itemize}}
\newcommand{\opt}[2]{\item[(#1)] #2}
% figure registry (generated by score2_tools.py finalize)
\qfigdef{p1_phy_s1_q01_1}{Score-II - Test 16 (09 Feb 2024)-img/p1_phy_s1_q01_1.png}{121.7}{0}
\qfigdef{p1_phy_s1_q02_1}{Score-II - Test 16 (09 Feb 2024)-img/p1_phy_s1_q02_1.png}{92.0}{0}
\qfigdef{p1_phy_s1_q04_1}{Score-II - Test 16 (09 Feb 2024)-img/p1_phy_s1_q04_1.png}{169.5}{0}
\qfigdef{p1_phy_s1_q05_1}{Score-II - Test 16 (09 Feb 2024)-img/p1_phy_s1_q05_1.png}{97.9}{0}
\qfigdef{p1_phy_s1_q06_1}{Score-II - Test 16 (09 Feb 2024)-img/p1_phy_s1_q06_1.png}{161.3}{0}
\qfigdef{p1_phy_s3_para1_2_1}{Score-II - Test 16 (09 Feb 2024)-img/p1_phy_s3_para1_2_1.png}{151.3}{0}
\qfigdef{p1_phy_s3_para3_4_1}{Score-II - Test 16 (09 Feb 2024)-img/p1_phy_s3_para3_4_1.png}{129.3}{0}
\qfigdef{p1_phy_s4_q08_1}{Score-II - Test 16 (09 Feb 2024)-img/p1_phy_s4_q08_1.png}{108.3}{0}
\qfigdef{p1_phy_s4_q08_2}{Score-II - Test 16 (09 Feb 2024)-img/p1_phy_s4_q08_2.png}{118.8}{0}
\qfigdef{p1_chem_s1_q01_1}{Score-II - Test 16 (09 Feb 2024)-img/p1_chem_s1_q01_1.png}{190.1}{0}
\qfigdef{p1_chem_s1_q03_A}{Score-II - Test 16 (09 Feb 2024)-img/p1_chem_s1_q03_A.png}{132.7}{0}
\qfigdef{p1_chem_s1_q03_B}{Score-II - Test 16 (09 Feb 2024)-img/p1_chem_s1_q03_B.png}{133.7}{0}
\qfigdef{p1_chem_s1_q03_C}{Score-II - Test 16 (09 Feb 2024)-img/p1_chem_s1_q03_C.png}{134.9}{0}
\qfigdef{p1_chem_s1_q03_D}{Score-II - Test 16 (09 Feb 2024)-img/p1_chem_s1_q03_D.png}{129.2}{0}
\qfigdef{p1_chem_s1_q04_1}{Score-II - Test 16 (09 Feb 2024)-img/p1_chem_s1_q04_1.png}{238.5}{0}
\qfigdef{p1_chem_s1_q04_A}{Score-II - Test 16 (09 Feb 2024)-img/p1_chem_s1_q04_A.png}{88.9}{0}
\qfigdef{p1_chem_s1_q04_B}{Score-II - Test 16 (09 Feb 2024)-img/p1_chem_s1_q04_B.png}{79.4}{0}
\qfigdef{p1_chem_s1_q04_C}{Score-II - Test 16 (09 Feb 2024)-img/p1_chem_s1_q04_C.png}{98.0}{0}
\qfigdef{p1_chem_s1_q04_D}{Score-II - Test 16 (09 Feb 2024)-img/p1_chem_s1_q04_D.png}{79.1}{0}
\qfigdef{p1_chem_s3_q02_1}{Score-II - Test 16 (09 Feb 2024)-img/p1_chem_s3_q02_1.png}{193.2}{0}
\qfigdef{p1_chem_s3_q04_1}{Score-II - Test 16 (09 Feb 2024)-img/p1_chem_s3_q04_1.png}{113.1}{0}
\qfigdef{p1_chem_s3_q04_2}{Score-II - Test 16 (09 Feb 2024)-img/p1_chem_s3_q04_2.png}{217.3}{0}
\qfigdef{p1_chem_s3_q05_1}{Score-II - Test 16 (09 Feb 2024)-img/p1_chem_s3_q05_1.png}{27.2}{0}
\qfigdef{p1_chem_s3_q05_2}{Score-II - Test 16 (09 Feb 2024)-img/p1_chem_s3_q05_2.png}{27.9}{0}
\qfigdef{p1_chem_s4_q07_1}{Score-II - Test 16 (09 Feb 2024)-img/p1_chem_s4_q07_1.png}{171.9}{0}
\qfigdef{p1_math_s2_para7_8_1}{Score-II - Test 16 (09 Feb 2024)-img/p1_math_s2_para7_8_1.png}{107.3}{0}
\begin{document}
\begin{center}{\Large\bfseries JEE (Main+Advanced) Score-II --- Test 16 (09 Feb 2024)}\end{center}
\begin{multicols}{2}
\partsec{1}{PHYSICS}

\sect{Section-I (i): One or More Correct}

\qn{1} Two identical massless rods of length 40 cm each are fixed to a rigid horizontal surface at points A \& B so that AB = 35 cm. Two identical point masses m = 1 kg each are permanently attached to the other ends of the rods as shown in the figure. The rods are connected to each other by a massless spring of force constant k = 100 N/m. The spring is attached to the point P \& Q of the rods where AP = BQ = 25 cm. Assume that the free length of the spring is 35 cm and it obeys Hooke's law. The rods are initially held at rest in a position when the angle $\theta$ made by each rod with the horizontal is equal to $37^\circ$. The rods are suddenly released from rest. The motion takes place in a vertical plane.
\qfig{p1_phy_s1_q01_1}
\begin{qoptions}
\opt{A}{The acceleration of point masses just after the system is released is approximately 23 $m/s^2$.}
\opt{B}{The angular velocity of rods when $\theta=90^\circ$ is approximately 8.36 rad/s.}
\opt{C}{The angular acceleration of rods when $\theta=90^\circ$ is zero.}
\opt{D}{The angular acceleration of rods when $\theta=90^\circ$ is 11.07 $rad/sec^2$.}
\end{qoptions}

\qn{2} Two moles of an ideal monoatomic gas is taken through a cycle ABCD as shown in the P-T diagram. During the process AB, pressure and temperature of the gas vary such that PT = constant. If $T_1=300$ K. (Take : $\ell n2=0.7$)
\qfig{p1_phy_s1_q02_1}
\begin{qoptions}
\opt{A}{The work done on the gas in the process AB is 1200 R Joules.}
\opt{B}{The work done on the gas in the process BC is $-600$ R Joules.}
\opt{C}{The work done by the gas in the process CA is 1200 R $\ell n2$ Joules.}
\opt{D}{The efficiency of the cyclic process is approximately 20\%.}
\end{qoptions}

\qn{3} A wire ACBOA shaped as a semi circle ACB with the bounding diameter BOA with OA = OB = a, has a total resistance R. It is placed in the same plane as infinitely long wire XY with AB parallel to the wire at a distance `a' and C further away from XY. A total charge Q flows through XY such that the current decreases linearly to zero from an initial value $I_0$ in time T.
\begin{qoptions}
\opt{A}{The magnetic flux through ACBOA at a time t is $\dfrac{\mu_0a}{\pi}\dfrac{Q}{T}(\pi-2)\left(1-\dfrac tT\right)$.}
\opt{B}{The induced current in ACBOA at a time t is $\dfrac{\mu_0a(\pi-2)}{\pi R}\dfrac{Q}{T^2}$.}
\opt{C}{The heat generated in ACBOA in time T is $\dfrac{a^2\mu_0^2(\pi-2)^2Q^2}{\pi^2RT^3}$.}
\opt{D}{The force acting on the ACBOA at a time t is $\dfrac{\mu_0^2Q^2(\pi-2)a}{2\pi^2RT^3}\left(1-\dfrac tT\right)$.}
\end{qoptions}

\qn{4} A disk of mass 4kg is mounted on a massless rigid rod BC as shown in the diagram. A torque $\tau=5e^{0.5t}$ is applied to the arm at C.
\qfig{p1_phy_s1_q04_1}
Consider three cases:\\
Case-I: The disk is free to rotate about its bearing at B.\\
Case-II: The disk is rigidly attached to its bearing at B.\\
Case-III: The disk is freely rotating about its bearing at B with angular velocity $(-10\hat k)$ rad/sec prior to application of torque. Take : e = 2.7
\begin{qoptions}
\opt{A}{In case-I, the kinetic energy of disk at t = 2 sec is $\dfrac{289}{2}$ J.}
\opt{B}{In case-II, the kinetic energy of disk at t = 2 sec is $\dfrac{1156}{9}$ J.}
\opt{C}{In case-III, the kinetic energy of disk at t = 2 sec is $\dfrac{339}{2}$ J.}
\opt{D}{In case-III, the kinetic energy of disk at t = 2 sec is $\dfrac{6281}{16}$ J.}
\end{qoptions}

\qn{5} A very large and thick sheet of conducting material carries a non-uniform current density $\vec J=J_0\left(\dfrac xd\right)^2\hat y$. The slab extends form x = $-$d to x = d, and is infinite in y and z.
\qfig{p1_phy_s1_q05_1}
\begin{qoptions}
\opt{A}{The magnetic field at $x=\dfrac d2$ is $\dfrac{\mu_0Jd}{24}$.}
\opt{B}{The magnetic field at x = 2d is $\dfrac{\mu_0Jd}{3}$.}
\opt{C}{The magnetic force per unit area on the right half of the sheet is $\dfrac{\mu_0J^2d^2}{18}$.}
\opt{D}{The magnetic force per unit area on the right half of the sheet is $\dfrac{\mu_0J^2d^2}{9}$.}
\end{qoptions}

\qn{6} An insulated small sphere carrying a positive charge is accelerated through a potential $V_1$, producing a horizontal beam along the x-axis. The beam enters a perfectly uniform electric field parallel to the y-axis, produced by a parallel plate capacitor with potential difference, $V_2$ between the plates, which have separation d and are of length, L as shown in figured. Take $V_1=2V$ and $V_2=1V$. Assume that collision of particle with the plates is perfectly elastic,
\qfig{p1_phy_s1_q06_1}
\begin{qoptions}
\opt{A}{If L = 100 d, the angle made by velocity of particle with the x-axis just after emerging from the capacitor is zero.}
\opt{B}{if L = 100 d, the number of collisions made by the particle with the plates of the capacitor is 25.}
\opt{C}{If L = 105d, the angle made by velocity of particle with the x-axis just after emerging from the capacitor is $\tan^{-1}\left(\dfrac12\right)$.}
\opt{D}{Irrespective of the value of L, the number of collisions between the particle and the upper plate of the capacitor is same.}
\end{qoptions}

\sect{Section-I (ii): Single Correct}

\paraq{7}{8}
Uncharged metallic spheres of radii 6R, 3R and 2R are mounted on insulated stands (separated by large distance). The spheres of radii 2R and 6R are charged to a potential V above earth potential. All three spheres are then briefly joined by a copper wire.

\qn{7} The final charge on sphere of radius 3R is :-
\begin{qoptions}
\opt{A}{$\dfrac{96\pi\varepsilon_0RV}{11}$}
\opt{B}{$\dfrac{192\pi\varepsilon_0RV}{11}$}
\opt{C}{$\dfrac{48\pi\varepsilon_0RV}{11}$}
\opt{D}{$\dfrac{8\pi\varepsilon_0RV}{11}$}
\end{qoptions}

\qn{8} The total heat loss in the process is:
\begin{qoptions}
\opt{A}{$\dfrac{128\pi\varepsilon_0V^2R}{11}$}
\opt{B}{$\dfrac{48\pi\varepsilon_0V^2R}{11}$}
\opt{C}{$\dfrac{96\pi\varepsilon_0V^2R}{11}$}
\opt{D}{$\dfrac{24\pi\varepsilon_0V^2R}{11}$}
\end{qoptions}

\paraq{9}{10}
When photons collide with hydrogen atoms in ground state, the atoms absorb the energy and rise to an excited state. After a time of $10^{-8}$ sec, these atoms release photons to come to the ground state. When the energy of a photon is greater than the ionization energy, it knocks out an electron from a atom. In a certain experiment, a stationary hydrogen atom is being bombarded with a source that a single photon strikes the atom after every few nanoseconds. The incident photons have different energies. The table below indicates the energies of the incident photon at different times. Assume hydrogen atom remains stationary during and after collision.
\begin{center}
\begin{tabular}{|c|c|c|}
\hline
Sl. No. & Time (in nanosecond) & Energy (eV) \\
\hline
1 & 100 & 10.2 \\
\hline
2 & 300 & 12.09 \\
\hline
3 & 400 & 10.2 \\
\hline
4 & 405 & 4.8 \\
\hline
5 & 550 & 12.75 \\
\hline
6 & 700 & 18 \\
\hline
\end{tabular}
\end{center}
Read the above passage carefully and answer the following questions.

\qn{9} After the first two photons are incident, the detector observes:
\begin{qoptions}
\opt{A}{one photon of energies 10.2 eV and an electron of energy 1.5 eV.}
\opt{B}{two photon of energies 10.2 eV and 12.09 eV.}
\opt{C}{two electrons of energies 3.4 eV and 1.51 eV.}
\opt{D}{no photon or electron emission.}
\end{qoptions}

\qn{10} Just after the 4th photon is incident, the detector observes:
\begin{qoptions}
\opt{A}{an electron having energy 4.8 eV.}
\opt{B}{a photon having energy 4.8 eV.}
\opt{C}{a photon having energy 8.8 eV.}
\opt{D}{an electron having energy 1.4 eV.}
\end{qoptions}

\sect{Section-II (i): Paragraph Type}

\paraq{1}{2}
The Bernoulli effect can have important consequences for the design of buildings. For example, wind can blow around a skyscraper at remarkably high speed, creating low pressure. The higher atmospheric pressure in the still air inside the buildings can cause windows to pop out. As originally constructed, the ALLEN Building popped windowpanes that fell from many stories to the sidewalk below.
\qfig{p1_phy_s3_para1_2_1}

\qn{1} Suppose a horizontal wind blows with a speed of 11.2 m/s outside a large pane of plate glass with dimensions 4.00 m $\times$ 1.50m. Assume the density of the air to be constant at 1.20 $kg/m^3$. The air inside the building is at atmospheric pressure. What is the total force exerted by air (in N) on the windowpane ?

\qn{2} If second Skyscraper is built near by, the air speed can be increased when wind passes through the narrow seperation between the building due to which force on the window pane having dimensions $4m\times1.5m$ becomes 1810 Newton outward. Find the velocity of the wind (in m/s) near the window pane after the construction of second building. Air density is 1.20 $kg/m^3$.

\paraq{3}{4}
The circuit in figure enables the measurement of the EMF, $\varepsilon$, of a thermocouple T (a temperature dependent source of a small EMF). AB is a uniform wire of length 1.00 m and resistance 2.00 $\Omega$. $K_1$ and $K_2$ are switches.\\
No current flows through G when:\\
(i) $K_1$ is closed and $K_2$ is open and AC = 90.0 cm.\\
(ii) $K_2$ is closed and $K_1$ is open and AC = 45.0 cm.
\qfig{p1_phy_s3_para3_4_1}

\qn{3} The EMF of the thermocouple (in mV) is:

\qn{4} The resistance R (in $\Omega$) is:

\paraq{5}{6}
In a YDSE experiment, the distance between the slits \& the screen is 100 cm . For a certain distance between the slits, an interference pattern is observed on the screen with the fringe width 0.25 mm. When the distance between the slits is increased by $\Delta d=1.2$ mm, the fringe width decreased to n = 2/3 of the original value. In the final position, a thin glass plate of refractive index 1.5 is kept in front of one of the slits \& the shift of central maximum is observed to be 20 fringe width.

\qn{5} The wavelength of light (in nm) is:

\qn{6} The thickness of the plate (in $\mu m$) is:

\sect{Section-II (ii): Numerical Value}

\qn{7} A solid square cross-section mild steel bar, of side 4cm is to be bent on the arc of a circle. The smallest radius (in m) to which it can be bent, if the breaking stress of the steel is 840 MPa, and Young's Modulus is 210 GPa is : (Assume that the radius of curvature is much larger than the thickness of the bar.)

\qn{8} A thin, convex lens $L_1$ forms a real image of an object located very far way, as shown in figure-(a). The image is located at a distance $4\ell$ and has height h. A concave lens $L_2$, of focal length $\ell$ is placed $2\ell$ from $L_1$. A convex lens, of focal length $2\ell$, is placed $3\ell$ from $L_1$, as shown in figure-(b). The height of final image is $\alpha h$, the value of $\alpha$ is :-
\qfig{p1_phy_s4_q08_1}
\qfig{p1_phy_s4_q08_2}

\qn{9} Living wood takes in radioactive carbon 14 from the atmosphere during the process of photosynthesis; the proportion of carbon 14 to carbon 12 being $1.25\times10^{-12}$. When the wood dies the carbon 14 decays, its half life being 5600 years. 4.00 g of carbon from a piece of dead wood gave a total count of 20.0 disintegrations per minute, with an uncertainty of 0.4 disintegrations per minute. The uncertainty in the age (in years) is. Take : $\ell n2=0.7$

\partsec{2}{CHEMISTRY}

\sect{Section-I (i): One or More Correct}

\qn{1} Consider the following system.\\
Three different aqueous solutions each having initial volume 100 ml are taken and kept in contact at constant temperature as shown in figure below :
\qfig{p1_chem_s1_q01_1}
After sufficient time (Consider 100\% dissociation of strong electrolytes)
\begin{qoptions}
\opt{A}{Decrease in volume of urea solution with respect to initial solution is $\dfrac{100}{3}$ml.}
\opt{B}{Final volume of $AlCl_3$ solution will be $\dfrac{400}{3}$ml.}
\opt{C}{There will be no change in volume of KCl solution.}
\opt{D}{Volume of both KCl and $AlCl_3$ solutions will increase.}
\end{qoptions}

\qn{2} Identify correct statement(s) from the following given options :
\begin{qoptions}
\opt{A}{The expression for \% ionisation of a dilute solution of a monoacidic weak base in aqueous solution is $\dfrac{100}{1+10^{(pK_b-pOH)}}$}
\opt{B}{If equal volumes of following solutions are mixed, then the pH of resulting solution will be average value of pH of two given solutions.\\
Solution I : pH = 5 ($CH_3COOH$)(aq) \& Solution-II : pH = 9 ($NH_3$(aq)),\\
Given : $[K_a(CH_3COOH)=K_b(NH_3(aq))]$}
\opt{C}{An acidic indicator HIn (colourless) will be visible only when its ionised form (pink in colour) is 25 \% by mole or more of total indicator in a solution. Suppose HIn ($pK_a=9.0$) is added to a solution having pH = 9.6 then mole \% of ionised form of HIn in solution is 90\%, (Take log 2 = 0.3)}
\opt{D}{Molar solubility of AgCl ($K_{SP}=10^{-12}$) in 0.2 M $MgCl_2$ aqueous solution is $2.5\times10^{-12}$ M.}
\end{qoptions}

\qn{3} Which of the following show \textbf{INCORRECT} major product ($E_2$ reaction)
\begin{qoptions}
\opt{A}{\qfigi{p1_chem_s1_q03_A}}
\opt{B}{\qfigi{p1_chem_s1_q03_B}}
\opt{C}{\qfigi{p1_chem_s1_q03_C}}
\opt{D}{\qfigi{p1_chem_s1_q03_D}}
\end{qoptions}

\qn{4} One or more than are correct.
\qfig{p1_chem_s1_q04_1}
Then the correct options is/are
\begin{qoptions}
\opt{A}{\qfigi{p1_chem_s1_q04_A}}
\opt{B}{\qfigi{p1_chem_s1_q04_B}}
\opt{C}{\qfigi{p1_chem_s1_q04_C}}
\opt{D}{\qfigi{p1_chem_s1_q04_D}}
\end{qoptions}

\qn{5} The \textbf{INCORRECT} statement(s) is/are :-
\begin{qoptions}
\opt{A}{The decomposition of ozone into oxygen is an exothermic process having negative Gibb's free energy change}
\opt{B}{Hydrogen peroxide is produced through catalytic hydrogenation of 2-Ethylanthraquinol}
\opt{C}{In froth flotation process non polar tails point outwards, making the detergent coated mineral particles hydrophilic}
\opt{D}{$Hg^{2+}$ ions give deep blue ppt. with cobalt thiocyanate}
\end{qoptions}

\qn{6} In which reaction(s) ; Phosphorus in reactant is trivalent but phosphorus is pentavalent in product.
\begin{qoptions}
\opt{A}{Reaction of phosphine with excess of nitrous oxide in presence of electrical sparking}
\opt{B}{Reaction of phosphorus oxychloride with excess amount of water}
\opt{C}{Reaction of phosphorus trichloride with acetic acid.}
\opt{D}{Reaction of white phosphorus with $SOCl_2$}
\end{qoptions}

\sect{Section-I (ii): Single Correct}

\paraq{7}{8}
For the first order dissociation reaction\\
$A_2B_3(aq.)\longrightarrow2A^{3+}(aq.)+3B^{2-}(aq.)$\\
a definite volume of solution is taken at different instants and its osmotic pressure rise is measured. If osmotic pressure at [t = 0] is 2 mm of Hg and at [t = 10] min is 6 mm of Hg then answer the questions given below for the solution :\\
\relax[Given density of solution remains constant and is equal to 1gm/ml]

\qn{7} Osmotic pressure of solution in mm of Hg at t = 20 min is :
\begin{qoptions}
\opt{A}{7 mm}
\opt{B}{8 mm}
\opt{C}{10 mm}
\opt{D}{12 mm}
\end{qoptions}

\qn{8} Time required for 87.5 \% of the reaction to complete is :
\begin{qoptions}
\opt{A}{15 min}
\opt{B}{20 min}
\opt{C}{30 min}
\opt{D}{40 min}
\end{qoptions}

\paraq{9}{10}
$X+O_2\xrightarrow{catalyst}Y(\text{neutral oxide})+Z(\text{neutral oxide})$\\
\resizebox{\linewidth}{!}{$Z\xrightarrow{electrolysis}A(\text{inflammable gas})+B(\text{supporter of combustion})$}\\
$Y+B\longrightarrow C(gas)$\\
$C(gas)+Z\longrightarrow E+D$\\
$E\xrightarrow{disproportionation}D+Y+B$

\qn{9} The sum of oxidation of central atom in X, Y and D is :-
\begin{qoptions}
\opt{A}{2}
\opt{B}{3}
\opt{C}{4}
\opt{D}{5}
\end{qoptions}

\qn{10} When excess of X is treated with aqueous solution of $CuSO_4$ then a complex is formed. Find the correct option regarding the complex ;
\begin{qoptions}
\opt{A}{Red color, paramagnetic, $\mu=1.73$ B.M.}
\opt{B}{Blue color, diamagnetic, $\mu=0$}
\opt{C}{Blue color, paramagnetic, $\mu=1.73$ B.M.}
\opt{D}{Red color, paramagnetic, $\mu=0$ B.M.}
\end{qoptions}

\sect{Section-II (i): Paragraph Type}

\paraq{1}{2}
An ideal mono-atomic gas is subjected to following series of steps.\\
Step-I : Isothermal single step expansion from (10 atm, 1 litre) to 1 atm at 300 K\\
Step-II : Isochoric heating to original pressure.\\
Step-III : Expansion from (10 atm., 10 litres) to (1 atm, 20 litres) in which pressure varies linearly with volume.\\
Step-IV : Pressure is again restored to original pressure at constant volume. (Use : 1 atm-litre = 100 J)

\qn{1} Magnitude of total work (in atm-litres) for the complete process is.

\qn{2} Change in internal energy (in kJ) for the complete process is :

\paraq{3}{4}
\qfig{p1_chem_s3_q02_1}

\qn{3} Number of correct statements in the following.\\
(a) (P) is phenol\\
(b) (S) is soluble in $NaHCO_3$\\
(c) (Z) is Cumene\\
(d) (Q) shows Iodo form test\\
(e) (Q) shows Tollen's test

\qn{4} 1 mole of (W) on reaction with $I_2$ in basic medium results in x gm of a yellow compound. Find out the value of x\\
(Atomic Wt. of I = 127)
\qfig{p1_chem_s3_q04_1}

\paraq{5}{6}
\qfig{p1_chem_s3_q04_2}

\qn{5} Which of the following statements is/are correct?\\
(a) (V) is \qfigi{p1_chem_s3_q05_1}\\
(b) (V) is \qfigi{p1_chem_s3_q05_2}\\
(c) (Q) is a $3^\circ$ -- Alcohol\\
(d) (U) is a $2^\circ$ -- Bromide\\
(e) (S) is $3^\circ$ -- Chloride\\
(f) Number of $\alpha$ -- H is (R) = 7

\qn{6} Mass of 1/9.6 mole of (V) is \_\_\_\_ gm.

\sect{Section-II (ii): Numerical Value}

\qn{7} \qfig{p1_chem_s4_q07_1}
Number of moles of Grignard reagent consumed per mole of above compound.

\qn{8} If an octahedral homoleptic complex $[Fe^{III}L_x]^{n+}$ does not exhibit geometrical isomerism but it exhibits optical isomerism, then find out the total number of correct statement(s) (denticity of L $\le$ 2).\\
(i) Number of ligands are 4\\
(ii) Number of S.I. possible are 2\\
(iii) Complex does not follow Sidgwick rule\\
(iv) Minimum and Max number of $e^-$ is present in $t_{2g}$ can be 3 and 5 respectively.\\
(v) The complex will always be paramagnetic in nature.\\
(vi) Complex will always be high spin complex

\qn{9} A black colored compound (A) on reaction with dilute $H_2SO_4$ gives a gas (B) which on passing in a solution of an acid (C) gives a white turbidity (D). Gas (B) when passed in an acidified solution of a compound (E) gives a precipitate (F) soluble in dilute $HNO_3$. After boiling this solution when an excess of $NH_4OH$ is added, a blue colored compound (G) is formed. To this solution on addition of acetic acid and aqueous $K_4[Fe(CN)_6]$ a chocolate precipitate (H) is obtained. On addition of an aqueous solution of $BaCl_2$ to an aqueous solution of (E), a white precipitate insoluble in $HNO_3$ is obtained.\\
Calculate the sum of oxidation state of every different type of metal present in the compounds from A to H.

\partsec{3}{MATHEMATICS}

\sect{Section-I (i): One or More Correct}

\qn{1} Let Set A = $\{1,2,3,4,\ldots,10\}$\\
and Set B = $\{2,3,4,\ldots,11\}$\\
then number of one-one functions defined from set A to set B such that
\begin{qoptions}
\opt{A}{$f(1)>f(2)>f(3)>f(4)<f(5)<f(6)<f(7)<f(8)<f(9)<f(10)$ is 84.}
\opt{B}{$f(x)\ne x+1$ for any x, is equal to 1334961}
\opt{C}{f is strictly increasing, is equal to 1}
\opt{D}{$f(x)\ne x+1$ for any x, is equal to $\dfrac{10!}{9}$}
\end{qoptions}

\qn{2} Let $z_1,z_2,z_3,\ldots,z_n$ be the vertices of a regular polygon and roots of the equation $z^n=2i$, if area of polygon is A, then correct option/s is/are
\begin{qoptions}
\opt{A}{$A=\dfrac n2\cdot2^{2/n}\cdot\sin\dfrac{2\pi}{n}$}
\opt{B}{$A=n\cdot2^{2/n}\cdot\sin\dfrac{2\pi}{n}$}
\opt{C}{$\displaystyle\lim_{n\to\infty}A=\pi$}
\opt{D}{$\displaystyle\lim_{n\to\infty}A=2\pi$}
\end{qoptions}

\qn{3} Let f(x) be a differentiable function satisfying $f(x+y)=5^xf(y)+3^yf(x)\ \forall\ x,y\in R$ but not always 0, then identify the correct option's
\begin{qoptions}
\opt{A}{$\dfrac{f(2)}{f'(2)}=\dfrac{16}{25\ell n5-9\ell n3}$}
\opt{B}{$f(0)=0$}
\opt{C}{$\dfrac{f(2)}{f(1)}=8$}
\opt{D}{if $f(1)=2$, then $f(-1)+f(-2)+f(-3)+\ldots\infty=-\dfrac14$}
\end{qoptions}

\qn{4} Let w be a variable complex number satisfying $|w-3-4i|=1$ and z be a complex number satisfying $amp\,\bar z=\tan^{-1}\dfrac43$ and $|z-3+2i|=2$ and $\alpha$ be a complex number in same argand plane, then correct option/s is/are
\begin{qoptions}
\opt{A}{$\left[\cos^{-1}\dfrac18,\cos^{-1}\left(-\dfrac18\right)\right]$}
\opt{B}{$\left[-\sin^{-1}\dfrac{\sqrt{35}}{6},\sin^{-1}\dfrac{\sqrt{35}}{6}\right]$}
\opt{C}{Minimum value of $|z-\alpha|+|w-\alpha|$ is 3}
\opt{D}{Minimum value of $|z-\alpha|+|w-\alpha|$ is 7}
\end{qoptions}

\qn{5} If $f(x)=\begin{cases}a[x]+bx^2+1, & -4\le x\le2\\x^2+a\{x\}+b, & 2<x\le5\\\dfrac{cx+4}{x-b}, & 5<x\le6\end{cases}$ is always continuous in it's domain then correct option's is/are\\
(where [.], \{.\} represents greatest integer function \& fractional part function respectively)
\begin{qoptions}
\opt{A}{$23b+a-c=3$}
\opt{B}{$\dfrac{3b-3a}{c-5b+a}=\dfrac15$}
\opt{C}{$\dfrac{c+2b+a}{c-2b-a}=\dfrac{11}{9}$}
\opt{D}{$\left[\dfrac{c}{10}\right]+\left[\dfrac b2\right]+\left[\dfrac a5\right]=2$}
\end{qoptions}

\qn{6} Let $z_1$ and $z_2$ be neither purely imaginary nor purely real such that $|z_1|=|z_2|=1$ (but $z_1\ne z_2$), if $\pm z_i,\pm\bar z_i\ (i\in\{1,2\})$ are the vertices of a convex polygon, z is a varibale complex number in same argand plane and $\lambda=|z-z_1|+|z-z_2|+|z+\bar z_2|+|z+\bar z_1|+|z+z_1|+|z+z_2|+|z-\bar z_2|+|z-\bar z_1|$ then correct option's is/are
\begin{qoptions}
\opt{A}{For $\lambda$ minimum, z must be O.}
\opt{B}{$\lambda_{min}=8$}
\opt{C}{If $amp\,z_1=\dfrac\pi8$ and $z_2=z_1\cdot e^{i\pi/4}$, then $\pm z_i,\pm\bar z_i$ (where $i\in\{1,2\}$) are $8^{th}$ roots of unity}
\opt{D}{If $amp\,z_1=\dfrac\pi8$ and $z_2=z_1\cdot e^{i\pi/4}$, then $\pm z_i,\pm\bar z_i\ (i\in\{1,2\})$ are the vertices of a 8 sided regular polygon}
\end{qoptions}

\sect{Section-I (ii): Single Correct}

\paraq{7}{8}
Consider a prism with triangular base $\triangle ABC$ such that equations of the rectangular surfaces ABB'A' and ACC'A' are $2x-y+z=2$ and $x+y+z=4$ respectively such that point C is $(0,0,4)$ and $AA'=6\sqrt{14}$ (as shown in fig.)
\qfig{p1_math_s2_para7_8_1}
On the basis of above details, answer the following

\qn{7} Volume of prism is
\begin{qoptions}
\opt{A}{$2\sqrt{14}$ cubic units}
\opt{B}{$\dfrac{6}{\sqrt{14}}$ cubic units}
\opt{C}{2 cubic units}
\opt{D}{4 cubic units}
\end{qoptions}

\qn{8} Equation of AC is
\begin{qoptions}
\opt{A}{$\dfrac x4=\dfrac y5=\dfrac{z-4}{-1}$}
\opt{B}{$\dfrac{x+4}{-4}=\dfrac{y-5}{5}=\dfrac{z-3}{-1}$}
\opt{C}{$\dfrac{x}{-4}=\dfrac y5=\dfrac{z-4}{1}$}
\opt{D}{$\dfrac{x-4}{-4}=\dfrac{y+5}{5}=\dfrac{z-3}{-1}$}
\end{qoptions}

\paraq{9}{10}
Variable circles $C_1,C_2,C_3,C_4$ with radii a,b,c,d respectively and centres above x-axis touch a fixed line y = 0 at the points $P(-4,0)$, $Q(4,0)$, R,S respectively. If circles $C_1$ and $C_2$ touches each other externally at M and $C_3$ touches $C_1,C_2$ and $C_4$ externally, also given that $C_4$ is externally tangent to $C_1$ On the basis of above answer the following.

\qn{9} Area bounded between locus of M and x-axis is
\begin{qoptions}
\opt{A}{$8\pi$}
\opt{B}{$10\pi$}
\opt{C}{$12\pi$}
\opt{D}{$\pi$}
\end{qoptions}

\qn{10} Correct option is
\begin{qoptions}
\opt{A}{$\sqrt a,\sqrt b,\sqrt c$ in Harmonic Progression}
\opt{B}{$\sqrt b,\sqrt c,\sqrt d$ in Harmonic Progression}
\opt{C}{a,b,c,d in Geometric Progression}
\opt{D}{a,b,c,d in Harmonic Progression}
\end{qoptions}

\sect{Section-II (i): Paragraph Type}

\paraq{1}{2}
Consider the function $f(x)=\dfrac{\displaystyle\sum_{r=1}^7(x-3r-1)^2}{282-\displaystyle\sum_{r=1}^{25}|x-r|}$

\qn{1} Minimum value of f(x) is equal to

\qn{2} Minimum value of f(x) occurs at x is equal to

\paraq{3}{4}
Tangents are drawn from any point lying on the ellipse $\dfrac{x^2}{16}+\dfrac{y^2}{4}=1$ to the circle $x^2+y^2=4$ such that the chord of contact always touches a fixed conic C. Let $F_1$ and $F_2$ be the foci of C and P be a point on C such that area of $\triangle PF_1F_2$ is maximum, then answer the following.

\qn{3} (Area of $\triangle PF_1F_2)^2$ is equal to

\qn{4} If a point moves inside the $\triangle PF_1F_2$ then the probability that it remains more close to side $F_1F_2$ as compared to sides $PF_1$ and $PF_2$ is equal to $\lambda$, then $\left(\dfrac{\lambda+3}{2}\right)^2$

\paraq{5}{6}
Let $f(x)=x^2-4x-4$, if a square is formed in the region bounded by $y=f(x)$ and x-axis such that one side of square on x-axis and other 2 vertices on $y=f(x)$, then

\qn{5} Area of the square is

\qn{6} $\displaystyle\int_0^4\dfrac{f(x)}{f(x)+f(4-x)}dx$ is equal to

\sect{Section-II (ii): Numerical Value}

\qn{7} There are 15 equispaced horizontal and 15 equispaced vertical lines such that distance between 2 consecutive lines is k units, let $S_i$ represents the sum of areas of all possible squares having side length (k.i), then $\dfrac{\displaystyle\sum_{i=1}^7S_i}{\displaystyle\sum_{i=8}^{14}S_i}$ is equal to

\qn{8} Let f(x) be a polynomial function satisfying $\displaystyle\int\dfrac{e^x(2x^2+2x^3-x+1)}{x^2+e^{2x}(2x^2-1)^2}dx=\tan^{-1}e^x\left(2f(x)-\dfrac{1}{f(x)}\right)+C$\\
then $f(2)-f(1)$ is equal to

\qn{9} If $A=\begin{pmatrix}1&-2\\0&1\end{pmatrix}$, $B=\begin{pmatrix}5&2\\2&1\end{pmatrix}$ and $C=A(ABA^T)^{25}(AB^TA^T)^{10}A^T$ then sum of elements of matrix C is

\end{multicols}
\end{document}
